\documentclass[10pt,a4paper]{amsart}
\usepackage[margin=19mm]{geometry}
\usepackage{amsmath,amssymb,mathtools}
\usepackage[scr=rsfs,cal=euler]{mathalfa}
\usepackage{xfrac}
\usepackage[unicode,hidelinks,bookmarks=false]{hyperref}
\newcommand{\ie}{i.e.\ }
\newcommand{\cf}{cf.\ }
\newcommand{\resp}{respectively\ }
\newcommand{\Subsection}{\subsection}
\providecommand{\nobreakdash}{\nobreak\hbox{-}}
\setlength{\emergencystretch}{2em}
\multlinegap0pt
\usepackage{xcolor}
\usepackage{amssymb}
\usepackage{tikz}
\usepackage{mathtools}
\usepackage[shortlabels]{enumitem}
\newtheorem{proposition}{Proposition}
\newtheorem{lemma}{Lemma}
\newtheorem{theorem}{Theorem}
\usepackage{cleveref}
\crefname{proposition}{Proposition}{Propositions}
\crefname{lemma}{Lemma}{Lemmas}
\crefname{theorem}{Theorem}{Theorems}
\newcommand\FF{\mathbb{F}}
\newcommand\ZZ{\mathbb{Z}}
\title{The Hasse norm principle for $A_4$-quartic extensions of global function fields}
\author{Anand Deopurkar, Rachel Newton, Vaidehee Thatte, Rosa Winter}
\date{Source: arXiv:2609.02444v1 (CC BY 4.0)}
\begin{document}
\maketitle

\noindent\emph{Source and licence.} The Hasse norm principle for $A_4$-quartic extensions of global function fields. Version arXiv:2609.02444v1 (CC BY 4.0), distributed by arXiv under the Creative Commons Attribution 4.0 International licence, http://creativecommons.org/licenses/by/4.0/, as recorded in the arXiv licence metadata for this exact version and shown on the abstract page https://arxiv.org/abs/2609.02444v1. Full licence notice: \texttt{licenses/arxiv-2609.02444-CC-BY-4.0.txt}.\par\smallskip

\noindent\textit{Editorial scope.} Counted unit: the article's theorem n (source0.tex lines 1297--1301). The article's complete original proof of it is delivered with it (source0.tex lines 1302--1357, one \texttt{proof} environment of 56 lines). No mathematics is written, repaired or re-derived: the statement and the proof are the article's own lines, in the article's own notation, and no retained line is substituted or re-flowed.  Retained uncounted as the source-local context the proof needs: lines 746--749; lines 1105--1114; lines 1220--1225; lines 1246--1248; lines 1259--1265.  Nothing was omitted from any retained span, so the package carries no omission record.  No retained line contains a citation, so the wrapper carries no bibliography and no bibliography entry.  No retained line contains a figure environment, an image-inclusion command or drawing code, so no image is delivered and the package declares no asset dependency.  Editorial formatting only, in addition: an AMS \texttt{amsart} preamble that restates the article's own declarations and macros used by the retained lines, namely the environments \texttt{proposition}, \texttt{lemma}, \texttt{theorem} on separate counters, together with the standard packages those lines need.  The retained environments print in this document as Proposition 1, Lemma 1, Theorem 1 in order of appearance, whereas the article prints the same items with the numbering of its own full document.  The complete native source of the article as distributed in the arXiv e-print for this exact version is bundled unchanged as \texttt{sources/209201/source0.tex}.
\begin{proposition}\label{prop:sterlingapprox}
  Let \(b_\ell = \sum_{i = 1}^\ell 2^{-i} s(\ell,i)/\ell!\).
  Then \(b_\ell < \ell^{-1/2}\) and \(b_\ell \sim \ell^{-1/2}\).
\end{proposition}

\begin{proposition}\label{prop:cpa4}
  \begin{enumerate}
  \item There exists $\lambda\in\mathbb{Z}_{>0}$ such that for all \(a,b,n_2,n \in \ZZ_{\geq 0}\) with $a+b+n_2=n$, 
    \[ |V(a,b,n_2)/{\mathbb{B}_n}| \leq \lambda.\]
   \item There exists $M\in\mathbb{Z}_{>0}$ such that for all \(a,b,n_2, n \in \ZZ_{\geq 0}\) with \(a,b,n_2 \geq M\) and $a+b+n_2=n$,
    \[ |V(a,b,n_2)/{\mathbb{B}_n}| = \begin{cases} 1 & \text{ if } a \equiv b \pmod 3 \\
    0 &\text{ otherwise.}
  \end{cases}\]
  \end{enumerate}
\end{proposition}

\begin{proposition}\label{prop:precise-q}
  Suppose \(a,b,n_2\) are such that \(V(\{a,b\},n_2)\) forms one orbit under \({\mathbb{B}_n} \times S_4\).
  Then \(H(\{a,b\}, n_2) \otimes \FF_p\) is geometrically connected and
  \[ \limsup_{m \to \infty} \frac{|\phi(H^{\rm fail}(\{a,b\},n_2)(\FF_{p^m}))|}{|H(\{a,b\},n_2)(\FF_{p^m})|} =\lim_{m \to \infty} \frac{|\phi(H^{\rm fail}(\{a,b\},n_2)(\FF_{p^{2m}}))|}{|H(\{a,b\},n_2)(\FF_{p^{2m}})|}= r b_{n_2},\]
   where $r\leq 2$ is the number of geometric components of \(\widetilde H^{\rm fail}(n_1,n_2) \otimes \FF_p\) lying above \(H(\{a,b\}, n_2) \otimes \FF_p\). 
\end{proposition}

\begin{equation}\label{eq:disjoint}
    H(n) =  \bigsqcup_{\substack{n_1+n_2 = n \\ n_1 > 0, n_2 \geq 0}} H(n_1,n_2).
\end{equation} 

\begin{lemma}\label{lem:estimates}
  Fix an integer \(m\).
  We have
  \[ \sum_{\substack{a+b+n_2 = n\\ a,b,n_2 \geq m,\ a \leq b}}v(\{a,b\},n_2) \approx n^2/12.\]
  Furthermore, for \(0 < \alpha < 1\), we have
  \[ \sum_{\substack{a+b+n_2 = n\\a,b,n_2 \geq m,\ a \leq b\\n_2 < n^{\alpha}}} v(\{a,b\},n_2) \approx n^{\alpha+1}/6.\]
\end{lemma}

\begin{theorem}\label{thm: failbound general n}
  Let \(p \geq 5\) be a prime.
    Given \(\epsilon > 0\), if \(n\) is sufficiently large, then 
  \[ \limsup_{m \to \infty} \frac{|\phi(H^{\rm fail}(n)(\FF_{p^m}))|}{|H(n)(\FF_{p^m})|} \leq (4+\epsilon) n^{-1/3}.\]
\end{theorem}

\begin{proof}
  All the summations that follow are over \(a,b,n_2\) satisfying \(a+b+n_2 = n\), $n_2\geq 0$ and \(a+b > 0\) and \(a \leq b\).
  Under the summation signs, we drop these conditions and only write additional conditions.
  Fix \(\alpha\) to be any real number in \((0,1)\).
  Let \(M\) be an integer such that \(v(\{a,b\},n_2) = 1\) for \(a,b,n_2 \geq M\) (see \Cref{prop:cpa4}).
  Using~\eqref{eq:disjoint} and
  \[H^{\rm fail}(n) = \bigsqcup_{\substack{a+b+n_2 = n\\ n_2\geq 0 \\ a+b > 0, a \leq b}} H^{\rm fail}(\{a,b\},n_2),\]
  when $n^\alpha\geq M$, we have
  \begin{equation}\label{eqn:3parts}
    \begin{split}
      |\phi(H^{\rm fail}(n)(\FF_{p^m}))| &= \sum_{\substack{a,b \geq M\\ n_2 \geq n^{\alpha}}}|\phi(H^{\rm fail}(\{a,b\},n_2)(\FF_{p^m}))| \\
                                        &+ \sum_{\substack{a, b, n_2 \geq M\\ n_2 < n^{\alpha}}} |\phi (H^{\rm fail}(\{a,b\},n_2)(\FF_{p^m}))| \\
                                        &+ \sum_{a,b,\text{ or } n_2 < M} |\phi (H^{\rm fail}(\{a,b\},n_2)(\FF_{p^m}))|.
    \end{split}
\end{equation}
We consider the three parts of the sum separately, in sequence.

  Let us consider the first part.
  If \(n_2 \geq n^{\alpha}\) and \(a,b \geq M\), \Cref{prop:precise-q} and \Cref{prop:sterlingapprox} imply 
  \[ \limsup_{m \to \infty} \frac{|\phi(H^{\rm fail}(\{a,b\},n_2)(\FF_{p^m}))|}{|H(\{a,b\},n_2)(\FF_{p^m})|}  < 2 n^{-\alpha/2} .\]
  Therefore,    \begin{equation}\label{eqn:part1}
    \limsup_{m \to \infty} \frac{\sum_{\substack{a,b \geq M\\ n_2 \geq n^{\alpha}}}|\phi(H^{\rm fail}(\{a,b\},n_2)(\FF_{p^m}))|}{|H(n)(\FF_{p^m})|}  < 2 n^{-\alpha/2} .
  \end{equation}

  Now let us consider the second and third parts.
  We have the bound
  \[|\phi (H^{\rm fail}(\{a,b\},n_2)(\FF_{p^m}))| \leq |H^{\rm fail}(\{a,b\},n_2)(\FF_{p^m})|.\]
  By the Lang--Weil bounds, we have
  \[ \limsup_{m \to \infty} \frac{|H^{\rm fail}(\{a,b\},n_2)(\FF_{p^m})|}{p^{mn}} \leq v(\{a,b\},n_2).\]
  Therefore, for $n$ sufficiently large, we have
\begin{equation}\label{eqn:part2}
    \begin{split}
      \limsup_{m \to \infty} &\frac{1}{p^{mn}}\left(\sum_{\substack{a, b, n_2 \geq M\\ n_2 < n^{\alpha}}} |\phi (H^{\rm fail}(\{a,b\},n_2)(\FF_{p^m}))| + \sum_{a,b,\text{ or } n_2 < M} |\phi (H^{\rm fail}(\{a,b\},n_2)(\FF_{p^m}))| \right) \\
      & \leq \sum_{\substack{a, b, n_2 \geq M\\ n_2 < n^{\alpha}}} v(\{a,b\},n_2) + \sum_{a,b,\text{ or } n_2 < M} v(\{a,b\},n_2) \approx n^{\alpha+1}/6.
    \end{split}
  \end{equation}

  
  The last approximation follows from \Cref{lem:estimates} applied to the first sum.
  By \Cref{prop:cpa4}, the second sum is much smaller---bounded by a multiple of \(n\).

  By Propositions~\ref{prop:cpa4} and~\ref{prop:precise-q}, the components \(H(\{a,b\},n_2) \otimes \FF_p\) of \(H(n) \otimes \FF_p\) for \(a,b,n_2 \geq M\) are geometrically connected, and again by the Lang--Weil bounds, we have
  \begin{equation}\label{eqn:part3}
    \liminf_{m \to \infty} \frac{1}{p^{mn}} |H(n)(\FF_{p^m})| \geq \sum_{a,b,n_2 \geq M} v(\{a,b\},n_2) \approx n^2/12.
  \end{equation}
  The last approximation is again thanks to \Cref{lem:estimates}.

  By combining the inequalities \eqref{eqn:part1}, \eqref{eqn:part2}, and \eqref{eqn:part3}, we get
  \begin{equation}
    \begin{split}
      \limsup_{m \to \infty} \frac{|\phi(H^{\rm fail}(n)(\FF_{p^m}))|}{|H(n)(\FF_{p^m})|} &< 2n^{-\alpha/2} + \frac{\sum_{\substack{a, b, n_2 \geq M\\ n_2 < n^{\alpha}}} v(\{a,b\},n_2) + \sum_{a,b,\text{ or } n_2 < M} v(\{a,b\},n_2)}{\sum_{a,b,n_2 \geq M}v(\{a,b\},n_2)}\\
      &\approx 2n^{-\alpha/2} + 2 n^{\alpha - 1}.
    \end{split}
  \end{equation}
Taking \(\alpha = 2/3\) yields the result.
\end{proof}

\end{document}
